\documentclass[11pt,a4paper]{article}

\usepackage[utf8]{inputenc}
\usepackage[T1]{fontenc}
\usepackage[spanish]{babel}
\usepackage[a4paper,margin=1.6cm]{geometry}
\usepackage{xcolor}
\usepackage{enumitem}
\usepackage{hyperref}
\usepackage{titlesec}
\usepackage{array}
\usepackage{tabularx}
\usepackage{lmodern}

\hypersetup{
    colorlinks=true,
    urlcolor=blue!50!black,
    linkcolor=black
}

\definecolor{primary}{RGB}{18, 52, 86}
\definecolor{accent}{RGB}{70, 70, 70}

\pagestyle{empty}
\setlength{\parindent}{0pt}
\setlength{\parskip}{0.35em}
\setlist[itemize]{leftmargin=1.2em, itemsep=0.2em, topsep=0.2em}

\titleformat{\section}
    {\large\bfseries\color{primary}}
    {}
    {0pt}
    {}
    [\titlerule]

\newcommand{\cvitem}[2]{
    \textbf{#1} \hfill {\color{accent} #2}
}

\newcommand{\smallsep}{\vspace{0.35em}}

\begin{document}

{\Huge \textbf{Juan Sebasti\'an Pe\~nalosa}}\\[0.2em]
{\Large Estudiante de F\'isica}\\[0.4em]
\href{https://www.linkedin.com/in/jsebastiands}{linkedin.com/in/jsebastiands}
\;|\;
\href{https://github.com/JsebastianUVPRQ}{github.com/JsebastianUVPRQ}
\;|\;
Ciudad: \textit{[Completar]}
\;|\;
Correo: \textit{[Completar]}

\section*{Perfil}
Estudiante de F\'isica con inter\'es en sistemas din\'amicos, estabilidad, caos, topolog\'ia,
aprendizaje autom\'atico y aplicaciones de la f\'isica cu\'antica a problemas reales.
Experiencia construyendo proyectos t\'ecnicos en Python, modelado cuantitativo,
an\'alisis de datos y prototipos de software para visualizaci\'on y predicci\'on.

\section*{Formaci\'on Acad\'emica}
\cvitem{Pregrado en F\'isica}{\textit{[Universidad]} -- \textit{[Fecha esperada de grado]}}\\
\'Enfasis de inter\'es: mec\'anica matem\'atica, sistemas din\'amicos, topolog\'ia,
f\'isica cu\'antica y modelado computacional.

\section*{Cursos y Certificaciones}
\cvitem{Advanced Linear Algebra}{Coursera}\\
Certificado verificable:
\href{https://coursera.org/verify/QLL6MYP6NF7T}{coursera.org/verify/QLL6MYP6NF7T}\\
Publicaci\'on relacionada:
\href{https://www.linkedin.com/posts/jsebastiands_completion-certificate-for-advanced-linear-share-7438070866820812801-Mcy_}{LinkedIn}

\smallsep
\cvitem{Estudio y divulgaci\'on en sistemas din\'amicos, estabilidad, caos y topolog\'ia}{Formaci\'on complementaria}\\
Evidencia p\'ublica:
\href{https://www.linkedin.com/posts/jsebastiands_sistemas-din%C3%A1micos-estabilidad-caos-y-topolog%C3%ADa-ugcPost-7397757554304643072-07vJ}{LinkedIn}

\smallsep
\cvitem{Exploraci\'on aplicada de f\'isica cu\'antica para detecci\'on de fraude}{L\'inea de estudio independiente}\\
Evidencia p\'ublica:
\href{https://www.linkedin.com/posts/jsebastiands_puede-la-f%C3%ADsica-cu%C3%A1ntica-detectar-fraudes-ugcPost-7417032038995886080-9PmG}{LinkedIn}

\section*{Habilidades T\'ecnicas}
\begin{itemize}
    \item \textbf{Lenguajes y herramientas:} Python, Jupyter Notebook, LaTeX, Git, GitHub.
    \item \textbf{An\'alisis y modelado:} \'algebra lineal, an\'alisis de datos, aprendizaje autom\'atico, modelado cuantitativo.
    \item \textbf{Bibliotecas y ecosistema:} NumPy, Pandas, SciPy, Folium, Streamlit.
    \item \textbf{Intereses de investigaci\'on:} informaci\'on cu\'antica, QML, sistemas no lineales, aplicaciones computacionales de la f\'isica.
\end{itemize}

\section*{Proyectos Seleccionados}
\cvitem{CaliCrashDashboard}{Python, Streamlit, Geoan\'alisis}\\
Dashboard para pron\'ostico y visualizaci\'on de siniestralidad vial en Cali, con enfoque en
datos georreferenciados, filtrado espacial y exploraci\'on interactiva.\\
Repositorio:
\href{https://github.com/JsebastianUVPRQ/CaliCrashDashboard}{github.com/JsebastianUVPRQ/CaliCrashDashboard}

\smallsep
\cvitem{QML}{Python, Quantum Machine Learning}\\
Proyecto orientado a modelos de aprendizaje autom\'atico inspirados en informaci\'on cu\'antica.\\
Repositorio:
\href{https://github.com/JsebastianUVPRQ/QML}{github.com/JsebastianUVPRQ/QML}

\smallsep
\cvitem{Machine\_learning\_project\_Flask\_deploy}{Python, Flask, ML}\\
Despliegue de un modelo de regresi\'on log\'istica mediante una aplicaci\'on web ligera.\\
Repositorio:
\href{https://github.com/JsebastianUVPRQ/Machine_learning_project_Flask_deploy}{github.com/JsebastianUVPRQ/Machine_learning_project_Flask_deploy}

\smallsep
\cvitem{ES\_Bogota\_real\_state\_market\_analysis}{An\'alisis de datos}\\
Exploraci\'on y estudio cuantitativo del mercado inmobiliario con notebooks reproducibles.\\
Repositorio:
\href{https://github.com/JsebastianUVPRQ/ES_Bogota_real_state_market_analysis}{github.com/JsebastianUVPRQ/ES_Bogota_real_state_market_analysis}

\section*{Divulgaci\'on e Intereses Acad\'emicos}
\begin{itemize}
    \item Inter\'es activo en divulgaci\'on cient\'ifica y matem\'atica.
    \item Curadur\'ia y publicaci\'on de contenido relacionado con f\'isica, matem\'aticas y tecnolog\'ia.
    \item Perfil de GitHub con proyectos en aprendizaje autom\'atico, f\'isica computacional y visualizaci\'on de datos.
\end{itemize}

\section*{Idiomas}
\begin{itemize}
    \item Espa\~nol: nativo.
    \item Ingl\'es: \textit{[Completar nivel: B2 / C1 / otro]}.
\end{itemize}

\section*{Notas para personalizaci\'on}
\begin{itemize}
    \item Reemplazar los campos \textit{[Completar]} por datos reales antes de enviar el CV.
    \item Si quieres una versi\'on orientada a investigaci\'on, se pueden a\~nadir secciones de seminarios, lecturas, p\'osters o inter\'eses de posgrado.
    \item Si quieres una versi\'on orientada a data/AI, conviene ampliar logros cuantificables por proyecto.
\end{itemize}

\end{document}
